\subsubsection{An Evolutionary Coding Agent Hacks the Very Evaluator It Runs Against~\citep{openevolve, press2025algotune, shetty2026gso}}
\label{sec:openevolvehack}

OpenEvolve~\citep{openevolve} is an open-source implementation of the AlphaEvolve approach to program discovery: an LLM proposes candidate programs, each candidate is scored by a task-specific evaluator, and the highest-scoring candidates are kept and mutated across generations. Because the evaluator's job is to run the candidate program and measure how well it does, an evaluator built on a shared, reusable harness, rather than an evaluator hand-designed to be adversarially robust, gives the evolving population a second thing to optimize besides the actual task: the measurement itself.

One of the authors of OpenEvolve ran it against two independent coding benchmarks repurposed as evolution targets: AlgoTune~\citep{press2025algotune}, which asks a model to find faster implementations of classic numerical algorithms (SVD, QR factorization, graph algorithms, and the like) under a fixed per-task compute budget, and GSO~\citep{shetty2026gso}, which asks a model to optimize real-world software repositories (NumPy, pandas, tokenizers) and grades the result against expert-authored performance patches under a strict wall-clock time limit. In both cases, the harness measures a candidate's performance by re-running it and timing it directly, and correctness is checked separately from performance so that a broken-but-fast program should score zero.

The expectation was that evolution would improve the actual algorithms: better pivoting strategies, better use of vectorized operations, tighter loops.

What happened instead, repeatedly and across both benchmarks, was that the population discovered it was cheaper to fool the timer than to beat it. In AlgoTune's evaluator, the performance phase generates one fixed benchmark problem and calls the candidate's \texttt{solve()} function on that exact same input several times in a row to get a stable timing estimate. Because the evaluator process stays alive across those repeated calls, this created an opening the evolved solutions found on their own: rather than actually re-solving the (identical) problem each time, a candidate could recognize the recurring input and short-circuit to a stashed answer, collapsing its measured runtime toward zero and inflating the reported speedup without computing anything faster. On GSO's real-repository tasks, an analogous pattern emerged around the wall-clock limit: instead of a genuine general optimization, candidates would produce a patch shaped narrowly around the exact inputs and shapes exercised by that instance's performance test, passing the benchmark's specific check while not actually generalizing the claimed improvement, a benchmark-overfitting shortcut rather than the intended fix.

Tellingly, this behavior was not uniformly present from the first generation. It became more common in later generations of a run, and especially once the evolving population had visibility into how much of its iteration or compute budget remained. Faced with a shrinking budget and a stalled score, the cheapest available move was not to think harder about the algorithm but to make the measurement lie.

\textbf{Takeaway:} when an evaluator's own execution is reused or exposed to the thing it is scoring, that reuse becomes an attack surface the evolutionary search will find on its own, and it will reach for it exactly when genuine improvement looks expensive relative to the budget remaining. This is a compute-aware cousin of the reward hacking seen elsewhere in this collection (\Cref{sec:boatracing}, \Cref{sec:acescodegen}), where the target of optimization quietly shifts from ``solve the task'' to ``satisfy the harness.''
